\begin{table}[htbp]
\centering
\caption{Comparación de estimaciones de transferencia CA bajo diferentes supuestos metodológicos. El modelo de referencia utiliza datos completos con previa Dirichlet uniforme. Las diferencias se expresan en puntos porcentuales respecto al baseline.}
\label{tab:sensitivity}
\small
\begin{tabular}{lccccc}
\toprule
\textbf{Prueba} & \textbf{CA$\to$FA (\%)} & \textbf{$\Delta$ (pp)} & \textbf{CA$\to$PN (\%)} & \textbf{$\Delta$ (pp)} & \textbf{N} \\
\midrule
Baseline (referencia)       & 46.5 & ---     & 53.5 & ---     & 7{,}271 \\
Previa conservadora         & 46.5 & 0.00    & 53.5 & 0.00    & 7{,}271 \\
Previa informativa          & 46.5 & 0.00    & 53.5 & 0.00    & 7{,}271 \\
Sin valores atípicos        & 47.8 & $+$1.31 & 52.2 & $-$1.31 & 7{,}163 \\
Bootstrap (media $\pm$ DE)  & 47.2 $\pm$ 3.0 & $+$0.70 & 52.8 $\pm$ 3.0 & $-$0.70 & 20 rep. \\
\bottomrule
\end{tabular}
\begin{flushleft}
\small \textit{Nota:} Las pruebas de sensibilidad previa comparan tres especificaciones Dirichlet ($\alpha = 0.5$, $1$, $2$). La prueba de valores atípicos excluye circuitos en el percentil 1\% superior e inferior de proporción CA. El bootstrap se realiza con 20 replicaciones (80\% de circuitos con reemplazo). DE: desviación estándar; pp: puntos porcentuales.
\end{flushleft}
\end{table}
